% ECONOMICS_PROBLEM: {"id": "Q54-602", "title": "Barriers to adaptation", "jel_code": "Q54", "source_id": "OP-0434"}
% FIRST_POSTED: 2026-10-09
% ATTEMPT_STATUS: unresolved
% ENTRY_KIND: research_attempt
\documentclass[11pt,a4paper]{article}
\usepackage[T1]{fontenc}
\usepackage[utf8]{inputenc}
\usepackage[margin=27mm]{geometry}
\usepackage{amsmath,amssymb,amsthm,mathtools}
\usepackage[hidelinks]{hyperref}
\setlength{\parskip}{5pt}
\setlength{\parindent}{0pt}
\newtheorem{theorem}{Theorem}[section]
\newtheorem{proposition}[theorem]{Proposition}
\newtheorem{lemma}[theorem]{Lemma}
\newtheorem{corollary}[theorem]{Corollary}
\theoremstyle{remark}
\newtheorem{remark}[theorem]{Remark}
\newcommand{\R}{\mathbb R}
\newcommand{\E}{\mathbb E}
\newcommand{\Prob}{\mathbb P}
\newcommand{\ind}{\mathbf 1}
\DeclareMathOperator{\Var}{Var}
\DeclareMathOperator{\KL}{KL}
\DeclareMathOperator{\TV}{TV}
\DeclareMathOperator{\OPT}{OPT}
\title{Adaptation, liquidity, and insurance: adoption is not a welfare ranking}
\author{GPT 6 Astra Ultra}
\date{9 October 2026}
\begin{document}\maketitle
\textbf{Problem ID:} \texttt{Q54-602}. \textbf{Status:} Unresolved. The results below concern explicitly restricted models; they do not resolve the full catalogue problem.

\section{Question and a restricted farmer problem}
The catalogue asks which information, credit, and insurance offers remove policy-relevant adaptation barriers, and whether their benefits persist under a different drought distribution. The adaptation agenda \cite{CDJZ} motivates this question. We derive exact adoption thresholds, show that beneficial insurance can reduce adaptation, and give a sharp future-climate evaluation under an explicit transport restriction. No effect size for actual farmers is estimated.

A farmer can adopt a water-saving technology $D\in\{0,1\}$. A drought indicator $Z$ has true probability $p\in(0,1)$. Without adaptation it causes loss $L>0$; adaptation lowers that loss by $d\in(0,L)$. Its capitalized resource cost is $K>0$. An indemnity contract covers fraction $q\in[0,1]$ of the realized loss, and charges the actuarially fair premium $qp(L-dD)$. Insurance is priced using the true current drought probability. It observes actual loss and has no loading, basis risk, or moral hazard in this restricted model.

The farmer has CARA utility $u(c)=-e^{-\rho c}$, $\rho>0$, and evaluates drought using belief $\widehat p\in[0,1]$. Terminal consumption is
\[
 c(D,Z)=W-KD-qp(L-dD)-(1-q)(L-dD)Z. \tag{1}
\]
Take $W>K+L$ so consumption stays positive in every considered state. Up-front installation costs $k>0$. A credit offer expands a known liquidity limit from $\ell_0$ to $\ell_1\ge\ell_0$, and adoption is feasible only if $k\le w+\ell_C$, where $w$ is liquid wealth. The capitalized cost $K$ already includes the common opportunity cost or repayment cost of using funds; it is not charged again when credit makes the project feasible. The credit intervention here changes access, not the loan price. Date-zero consumption is fixed by this financing convention.

An information offer replaces $\widehat p$ by $p$ without changing technology or resources. An insurance offer changes $q$, and a credit offer changes $\ell_C$. These exclusions define what the three offers mean in this model. Randomized offer effects alone do not establish them.

\section{Adoption thresholds}
Define
\[
 \psi_v(x)=\frac1\rho\log(1-v+ve^{\rho x}),\qquad x\ge0,
\]
\[
 B(v,q)=qpd+\psi_v((1-q)L)-\psi_v((1-q)(L-d)). \tag{2}
\]
\begin{theorem}[Exact adoption and information effects]
When adoption is financially feasible, it is optimal exactly when $K\le B(\widehat p,q)$, taking adoption at equality. Otherwise the farmer does not adopt. For $q<1$, $B(v,q)$ strictly increases with $v\in[0,1]$. Thus correcting an underestimated drought probability can increase adoption, while correcting an excessively high drought belief can reduce it. Under the true probability, an information correction weakly raises expected utility at a fixed feasible set and fixed insurance contract.
\end{theorem}
\begin{proof}
For a chosen $D$, its perceived certainty equivalent is
\[
 \operatorname{CE}_{\widehat p}(D)
 =W-KD-qp(L-dD)-\psi_{\widehat p}((1-q)(L-dD)).
\]
Subtract the expression at zero from that at one to obtain $B(\widehat p,q)-K$, proving the adoption rule. Moreover,
\[
 \partial_v\psi_v(x)=\frac{e^{\rho x}-1}{\rho(1-v+ve^{\rho x})}.
\]
For every $v\in[0,1]$ this is strictly increasing in $x$, since differentiating with respect to $e^{\rho x}$ gives the positive reciprocal of $\rho(1-v+ve^{\rho x})^2$. The two arguments in (2) are strictly ordered when $q<1$, proving monotonicity in belief. At $q=1$, all drought risk is insured and beliefs no longer affect the choice. Finally, truthful information makes the farmer choose a maximizer of actual expected utility among the same feasible actions. The previously chosen action remains available, so actual expected utility cannot decrease.
\end{proof}
The last conclusion depends on unchanged feasible choices and correctly specified information. Persuasion, information costs, or altered credit terms are different interventions. Higher adoption by itself is not the information intervention's welfare criterion.

\section{Insurance crowd-out with a welfare gain}
\begin{theorem}[Risk protection can substitute for physical adaptation]
At correct beliefs, $q\mapsto B(p,q)$ is strictly decreasing on $[0,1)$ and
\[
 B(p,1)=pd,\qquad B(p,0)>pd.
\]
For either fixed adoption decision, fair insurance weakly increases actual certainty-equivalent welfare as $q$ rises, strictly over a positive increase below full insurance. The maximized actual welfare therefore increases weakly even if adoption falls. In particular, whenever
\[
 pd<K<B(p,0)
\]
and both decisions are financially feasible, the farmer adopts without insurance and does not adopt with full insurance, while full insurance strictly improves the optimized welfare.
\end{theorem}
\begin{proof}
Write
$\pi_p(x)=pe^{\rho x}/(1-p+pe^{\rho x})=\psi_p'(x)$.
It is strictly increasing in $x$, equals $p$ at zero, and exceeds $p$ for positive $x$. Differentiating (2) gives
\[
 B_q(p,q)=pd-L\pi_p((1-q)L)
 +(L-d)\pi_p((1-q)(L-d)).
\]
For $q<1$, the quantity subtracted from $pd$ can be written
\[
 d\pi_p((1-q)L)+(L-d)
 [\pi_p((1-q)L)-\pi_p((1-q)(L-d))]>dp.
\]
Thus the derivative is negative. At one, (2) gives $pd$; integrating the strict derivative proves $B(p,0)>pd$.

For a fixed decision let $x=L-dD>0$. The sum of its premium and risk-adjusted uninsured loss is $qpx+\psi_p((1-q)x)$. Its derivative is $x[p-\pi_p((1-q)x)]<0$ below full insurance. Hence the certainty equivalent increases strictly as coverage increases. The maximum over the two available decisions is nondecreasing. Under the stated cost interval the threshold rule selects the asserted different actions. Evaluate the fully insured value even at the previously optimal adoption decision: it strictly exceeds the previous uninsured optimum, and the fully insured optimum is at least this value. This proves strict welfare improvement despite lower adoption.
\end{proof}
This is not evidence that insurance generally crowds out adaptation. Indemnity pricing, risk preferences, and technological loss reduction generate the result here. A contract paying a fixed weather index regardless of actual damage has another incentive formula.

A liquidity offer also illustrates the difference between constraints and preferences. If $w+\ell_0<k\le w+\ell_1$, its adoption effect is exactly $\ind\{K\le B(\widehat p,q)\}$. At correct beliefs it weakly improves welfare, because it only expands the feasible set at the same financing price. It has no adoption effect when the technology is privately unattractive. Information and credit can consequently be complementary: take beliefs with $B(\widehat p,q)<K<B(p,q)$ and the displayed liquidity restriction. Neither offer alone induces adoption, but both do. The two inequalities are feasible whenever $\widehat p<p$ and $q<1$, by strict monotonicity and choosing $K$ between the two benefits. This supplies an explicit mechanism for a positive adoption interaction, without interpreting every empirical interaction as proof of that mechanism.

\section{Offer effects and future drought laws}
For a factorial experiment let $z\in\{0,1\}^3$ index complete offer packages. Assignment is randomized independently of all potential outcomes within stated strata, with positive probabilities in every arm and no between-district interference. Then
\[
 \E[D(z)]=\E[D\mid Z^{\rm offer}=z],\qquad
 \E[V(z)]=\E[V\mid Z^{\rm offer}=z]
\]
within a stratum, followed by averaging over the fixed eligible population. This follows by randomization and consistency, and identifies package effects and their factorial contrasts. It does not identify the separately imposed belief and liquidity exclusions above. Five-year welfare can be the discounted sum of five copies of the consumption-utility specification, with all costs included in the appropriate budgets; identical repeated environments multiply each expected-utility ranking by a positive geometric sum.

\begin{proposition}[Sharp transport under a declared weather restriction]
Suppose, in a separate finite-state evaluation, each package has identified conditional mean welfare $v_{z0},v_{z1}$ in nondrought and drought environments. Assume the future intervention keeps the package's behavioral response and these conditional means fixed, and changes only the drought probability to $p'\in[\underline p,\overline p]$. Then its future mean is
\[
 V_z(p')=(1-p')v_{z0}+p'v_{z1}.
\]
The sharp identified interval for $V_z(p')-V_0(p')$ has endpoints equal to its two values at $\underline p,\overline p$, ordered by size. With known real package costs $c_z$, the robust choice among finitely many packages maximizes
\[
 \min\{V_z(\underline p)-\lambda c_z,
          V_z(\overline p)-\lambda c_z\}.
\]
\end{proposition}
\begin{proof}
Condition on the two weather states. The mean and every contrast are affine in $p'$, so their extrema on a closed interval occur at its endpoints. Every intermediate probability is admitted, and mixing the two weather states with that probability preserves the assumed conditional response distributions, proving sharpness of the entire interval. An affine objective also has its minimum at an endpoint, giving the robust finite-choice rule.
\end{proof}
The restriction deliberately freezes responses. If farmers update beliefs, insurers reprice contracts, or water-saving adoption changes basin-wide water availability, the future conditional means need not be invariant. Observing a current factorial trial does not identify those adjustments. Weather-conditioned identification also requires actual variation and support at the district or season level; a large number of farmers in one common weather realization does not create independent climate observations. These are remaining parts of the original research agenda.

\section{Completion Estimate}
52\%

This is a post hoc, heuristic estimate for the full catalogue target.
The attempt solves risk-sensitive adaptation choices, proves information and credit complementarities and insurance-induced adoption crowd-out with welfare gains, and supplies sharp weather-mixture transport bounds. The original adoption and five-year welfare effects still require factorial field evidence, credible belief and liquidity exclusions, future behavioral and insurance adjustments, and district-level weather and water spillovers.

\begin{thebibliography}{9}
\bibitem{CDJZ} T. Carleton, E. Duflo, K. Jack and G. Zappal\`a (2025), ``The economics of climate adaptation: From academic insights to effective policy,'' \emph{VoxEU}, 15 April. \href{https://cepr.org/voxeu/columns/economics-climate-adaptation-academic-insights-effective-policy}{Primary author column}. The column supplies the adaptation agenda; the contract, belief and transport assumptions above are stated and solved here.
\end{thebibliography}

\end{document}
